\documentclass[10pt]{article}

\usepackage[a4paper, total={6in, 9.5in}]{geometry}
\usepackage[dvipsnames,svgnames]{xcolor}
\usepackage[shortlabels]{enumitem}
\usepackage[framemethod=TikZ]{mdframed}
\usepackage{amsmath,amssymb,amsthm}
\usepackage[colorlinks]{hyperref}
\usepackage{mathrsfs}
\usepackage{tikz-cd}

\hypersetup{
	colorlinks=true,
	linkcolor=blue,
	filecolor=blue,      
	urlcolor=blue,
	pdftitle={MAT535 Notes}
}

\usepackage{mathtools}
\usepackage{thmtools}
\usepackage[textsize=scriptsize]{todonotes}

\renewcommand{\tilde}{\widetilde}
\renewcommand{\hat}{\widehat}
\renewcommand{\setminus}{\!-\!}
\providecommand{\ol}{\overline}
\providecommand{\eps}{\varepsilon}
\providecommand{\half}{\frac{1}{2}}
\providecommand{\dang}{\measuredangle} %% Directed angle
\providecommand{\dd}{\mathrm d}
\providecommand{\CC}{\mathbb C}
\providecommand{\FF}{\mathbb F}
\providecommand{\HH}{\mathbb H}
\providecommand{\NN}{\mathbb N}
\providecommand{\QQ}{\mathbb Q}
\providecommand{\RR}{\mathbb R}
\providecommand{\ZZ}{\mathbb Z}
\providecommand{\RP}{\mathbb RP}
\providecommand{\CP}{\mathbb CP}
\providecommand{\dg}{^\circ}
\providecommand{\ind}[1]{\mathbf 1_{#1}}
\providecommand{\ii}{\item}
\providecommand{\alert}[1]{{\sffamily\textbf{\textcolor{magenta}{#1}}}}
\providecommand{\opname}{\operatorname}
\providecommand{\li}{\opname{li}}
\providecommand{\ls}{\opname{ls}}
\providecommand{\supp}{\opname{supp}}
\providecommand{\aut}{\opname{Aut}}
\providecommand{\inv}{^{-1}}
\providecommand{\setdiff}{\smallsetminus}
\providecommand{\mb}[1]{\mathbf{#1}}
\providecommand{\abs}[1]{\left\lvert{#1}\right\rvert}
\providecommand{\norm}[1]{\left\|{#1}\right\|}
\providecommand{\dif}{\;d}

\reversemarginpar
\setlength{\marginparwidth}{1in}

\mdfdefinestyle{mdbluebox}{roundcorner=10pt,innerbottommargin=9pt,
	linecolor=blue,backgroundcolor=TealBlue!5,}
\declaretheoremstyle[headfont=\sffamily\bfseries\color{MidnightBlue},
mdframed={style=mdbluebox},]{thmbluebox}

\mdfdefinestyle{mdredbox}{frametitlefont=\bfseries,innerbottommargin=8pt,
	nobreak=true,backgroundcolor=Salmon!5,linecolor=RawSienna,}
\declaretheoremstyle[headfont=\bfseries\color{RawSienna},
mdframed={style=mdredbox},headpunct={\\[3pt]},postheadspace=0pt,]{thmredbox}

\mdfdefinestyle{mdgreenbox}{linecolor=ForestGreen,backgroundcolor=ForestGreen!5,
	linewidth=2pt,rightline=false,leftline=true,topline=false,bottomline=false,}
\declaretheoremstyle[headfont=\bfseries\sffamily\color{ForestGreen!70!black},
mdframed={style=mdgreenbox},headpunct={ --- },]{thmgreenbox}

\mdfdefinestyle{mdorangebox}{linecolor=RedOrange,backgroundcolor=RedOrange!5,
	linewidth=2pt,rightline=false,leftline=true,topline=false,bottomline=false,}
\declaretheoremstyle[headfont=\bfseries\sffamily\color{RedOrange!70!black},
mdframed={style=mdorangebox},headpunct={ --- },]{thmorangebox}

\mdfdefinestyle{mdblackbox}{linecolor=black,backgroundcolor=RedViolet!5!gray!5,
	linewidth=3pt,rightline=false,leftline=true,topline=false,bottomline=false,}
\declaretheoremstyle[mdframed={style=mdblackbox}]{thmblackbox}

\declaretheoremstyle[spaceabove=6pt,spacebelow=6pt]{boldhead}

\declaretheorem[style=thmbluebox,name=Theorem,numberwithin=section]{theorem}

\declaretheorem[style=boldhead,name=Example,sibling=theorem]{example}

\declaretheorem[style=thmbluebox,name=Corollary,sibling=theorem]{corollary}
\declaretheorem[style=thmbluebox,name=Proposition,sibling=theorem]{prop}
\declaretheorem[style=thmbluebox,name=Lemma,sibling=theorem]{lemma}
\declaretheorem[style=thmbluebox,name=Theorem,numbered=no]{theorem*}
\declaretheorem[style=thmbluebox,name=Lemma,numbered=no]{lemma*}
\declaretheorem[style=thmgreenbox,name=Claim,numberwithin=theorem]{claim}
\declaretheorem[style=thmgreenbox,name=Claim,numbered=no]{claim*}
\declaretheorem[style=boldhead,name=Remark,sibling=theorem]{remark}
\declaretheorem[style=boldhead,name=Remark,numbered=no]{remark*}
\declaretheorem[style=thmgreenbox,name=Definition,sibling=theorem]{definition}
\declaretheorem[style=thmgreenbox,name=Definition,numbered=no]{definition*}
\declaretheorem[style=thmorangebox,name=Warning,sibling=theorem]{warning}
\declaretheorem[style=thmorangebox,name=Warning,numbered=no]{warning*}
\declaretheorem[style=thmorangebox,name=Note,numbered=no]{note}
\declaretheorem[style=thmblackbox,name=Exercise,sibling=theorem]{exercise}
\declaretheorem[style=thmblackbox,name=Exercise,numbered=no]{exercise*}

\newenvironment{soln}{\begin{proof}[Solution]}{\end{proof}}
\newenvironment{walkthrough}{\noindent \textbf{Walkthrough.}}{\medskip}
\newlist{walk}{enumerate}{3}
\setlist[walk]{label=\bfseries (\alph*)}

\providecommand{\pow}{\operatorname{Pow}}
%\providecommand{\vocab}[1]{{{\sffamily\textolor{ForestGreen}{#1}}}}
\providecommand{\vocab}[1]{{{\sffamily\textit{#1}}}}
\providecommand{\lcm}{\operatorname{lcm}}
\providecommand{\ord}{\operatorname{ord}}
\providecommand{\Int}{\operatorname{Int}}
\providecommand{\rk}{\operatorname{rk}}
\providecommand{\Mat}{\operatorname{Mat}}

\usepackage{asymptote}
\begin{asydef}
	defaultpen(fontsize(10pt));
	unitsize(1cm);
	size(8cm); // set a reasonable default
	usepackage("amsmath");
	usepackage("amssymb");
	settings.tex="pdflatex";
	settings.outformat="pdf";
	// Replacement for olympiad+cse5 which is not standard
	import geometry;
	// recalibrate fill and filldraw for conics
	void filldraw(picture pic = currentpicture, conic g, pen fillpen=defaultpen, pen drawpen=defaultpen)
	{ filldraw(pic, (path) g, fillpen, drawpen); }
	void fill(picture pic = currentpicture, conic g, pen p=defaultpen)
	{ filldraw(pic, (path) g, p); }
	// some geometry
	pair foot(pair P, pair A, pair B) { return foot(triangle(A,B,P).VC); }
	pair orthocenter(pair A, pair B, pair C) { return orthocentercenter(A,B,C); }
	pair centroid(pair A, pair B, pair C) { return (A+B+C)/3; }
	// cse5 abbreviations
	path CP(pair P, pair A) { return circle(P, abs(A-P)); }
	path CR(pair P, real r) { return circle(P, r); }
	pair IP(path p, path q) { return intersectionpoints(p,q)[0]; }
	pair OP(path p, path q) { return intersectionpoints(p,q)[1]; }
	path Line(pair A, pair B, real a=0.6, real b=a) { return (a*(A-B)+A)--(b*(B-A)+B); }
	// cse5 more useful functions
	picture CC() {
		picture p=rotate(0)*currentpicture;
		currentpicture.erase();
		return p;
	}
	pair MP(Label s, pair A, pair B = plain.S, pen p = defaultpen) {
		Label L = s;
		L.s = "$"+s.s+"$";
		label(L, A, B, p);
		return A;
	}
	pair Drawing(Label s = "", pair A, pair B = plain.S, pen p = defaultpen) {
		dot(MP(s, A, B, p), p);
		return A;
	}
	path Drawing(path g, pen p = defaultpen, arrowbar ar = None) {
		draw(g, p, ar);
		return g;
	}
\end{asydef}

\newcommand{\scr}[1]{\mathscr{#1}}
\newcommand{\obj}{\opname{obj}}
\newcommand{\Ch}{\opname{Ch}}
\renewcommand{\hom}{\opname{Hom}}
\newcommand{\id}{\opname{id}}
\newcommand{\im}{\opname{im}}
\newcommand{\coker}{\opname{coker}}
\newcommand{\tensor}{\otimes}
\newcommand{\dirsum}{\oplus}
\newcommand{\bigtensor}{\bigotimes}
\newcommand{\bigdirsum}{\bigoplus}


\title{\vspace{-2.0cm}Fields and Galois Theory}
\author{\large Aditya Pahuja}
\date{\large \today}
\begin{document}
\maketitle
\tableofcontents
\pagebreak

Notes from D\&F Chapters 13 and 14.

\section{Basic definitions}
\subsection{Fields}
Recall that a \vocab{field} is a nonzero integral domain (so $0\ne 1$)
in which every nonzero element is invertible;
equivalently it is a commutative ring with no nontrivial ideals.
The latter description of fields implies that
a homomorphism of fields must be injective,
since $0\ne 1$ implies that the kernel cannot be the entire field.

\begin{definition}[Characteristic of a field]
    \label{def:characteristic}
    The \vocab{characteristic} of a field $F$
    is the additive order of $1$ in the field
    (i.e. the order of $1$ in $(F,+)$),
    or $0$ if $1$ does not have finite additive order.
\end{definition}

\begin{exercise}
    \label{exer:characteristic-is-prime}
    Show that the characteristic of a field is prime.
\end{exercise}

\todo{prime subfield}

\subsection{Field extensions}
The reason historically for thinking about Galois theory
was to talk about solutions to polynomial equations
with coefficients in a particular field.
The central object in this way of thinking is that of a field extension:
\begin{definition}[Field extension]
    \label{def:field-extension}
    If $F$ is a field and $K$ is a field containing $F$,
    we write that $K/F$ (pronounced ``$K$ over $F$'')
    is a \vocab{(field) extension} of $F$.
    One can also think of this as a homomorphism $F\to K$,
    since field homomorphisms are always injective.
\end{definition}

\begin{example}
    \label{ex:extensions}
    (1) Every field is an extension of its prime subfield.

    (2) The polynomial $x^2 + 1$ is irreducible in $\QQ[x]$,
    so $\QQ[x]/(x^2 + 1)$ is a field containing $\QQ$,
    hence is an extension of $\QQ$.
    This is isomorphic to the subfield $\QQ(i)\subseteq\CC$
    generated by $\QQ$ and $i$ (or a definition of the subfield,
    if you define $\CC = \RR[x]/(x^2 + 1)$).
\end{example}

Indeed, quotients by polynomials in this way
capture the idea of ``adjoining a root of $p$ to the field,''
which is why field extensions contain the relevant algebraic data
for trying to solve polynomials. That is,
if $F$ is a field and $p(x)\in F[x]$ an irreducible polynomial,
we can construct the extension $F[x]/(p(x))$ of $F$
in which $p$ necessarily has a root $\ol x$
(the image of $x$ in the quotient).

If $K/F$ is an extension of $F$, then $K$ must be closed under addition
and under multiplication by elements of $F$,
so it is naturally an $F$-vector space.
\begin{definition}[Degree of an extension]
    \label{def:degree}
    The \vocab{degree} of the extension $K/F$
    is the dimension of $K$ as an $F$-vector space,
    and it is denoted by $[K:F]$.
    If $[K:F]$ is finite we say that the extension is \vocab{finite}.
\end{definition}
The similarity in notation to the notation for the index of a subgroup
is not an accident; we will see in fact that this notation
is extremely fitting.

The other notational similarity in the definition of the degree
is the choice of word: this is not a coincidence either.
\begin{prop}
    \label{prop:degree-of-extension-by-poly}
    Let $F$ be a field,
    $p(x)\in F[x]$ an irreducible polynomial of degree $n$,
    and $K = F[x]/(p(x))$. Then $[K:F] = n$
    with $F$-basis $\ol 1$, $\ol x$, \dots, $\ol x^{n-1}$.
\end{prop}

This proposition follows from the Euclidean algorithm,
where the image of a polynomial in $F[x]$ ends up just being
its remainder upon division by $p(x)$.

\todo{write out an example}

Note that the Euclidean algorithm also gives you a procedure
for computing the inverse of an element in $K$.
\begin{exercise}
    \label{exer:inverse-of-x}
    What is the inverse of $\ol x$ in $\QQ[x]/(p(x))$?
    (You should be able to express your answer nicely
    in terms of $p(x)$.)
\end{exercise}

\begin{example}
    \label{ex:Q-cbrt-3}
    A very, very important counterexample in Galois theory
    is the field extension $\QQ(\sqrt[3]2) \cong \QQ[x]/(x^3 - 2)$.
    Observe that this field does \emph{not} contain
    all the roots of $x^3 - 2$.
\end{example}

Very often we want to think of an extension of $F$
as being generated by $F$ plus a couple of other elements
from some bigger ambient field (e.g. $\QQ(i)$ from earlier,
or $\QQ[x]/(x^2 - 2)\cong \QQ(\sqrt 2)\subseteq\CC$).
Supposing that $K/F$ is an extension of $F$
and $(\alpha_i)_{i\in I}$ are some elements in $K$,
we write $F((\alpha_i)_{i\in I})\subseteq K$ to be the extension of $F$
generated by the $\alpha_i$; often the $\alpha_i$ are finite in number,
in which case we write $F(\alpha_1,\dots,\alpha_n)$
and say that the extension is finitely generated.
If a field $K/F$ can be written as $F(\alpha)$ for some $\alpha\in K$
then we say that $K$ is a \vocab{simple extension} of $F$
and that $\alpha$ is a \vocab{primitive element} of the extension.

\begin{prop}
    \label{prop:simple-extension-quotient-poly}
    If the primitive element $\alpha$ of a simple extension $F(\alpha)$
    is the root of a polynomial $p(x)\in F[x]$
    which is irreducible over $F$, then
    \begin{equation*}
	F(\alpha)\cong F[x]/(p(x)).
    \end{equation*}
\end{prop}

The proof is a straightforward application of the first isomorphism theorem,
as we have a natural homomorphism $F[x]\to F(\alpha)$
by sending $x$ to $\alpha$, whose kernel can be shown to be $(p(x))$.
For an irreducible polynomial $p(x)\in F[x]$,
if $E/F$ is an extension containing a root $\alpha$ of $p(x)$
and $E = F(\alpha)$, we will call $E$ a \vocab{stem field} of $p(x)$.
This should be thought of as a smallest extension of $F$
in which $p$ has a root.

\begin{example}
    \label{ex:complex-cbrt-3}
    If $p(x) = x^3 - 2$ and
    $\alpha = \sqrt[3]2i$, then $\QQ(\alpha)\subseteq\CC$
    is isomorphic to $\QQ[x]/(x^3 - 2)$,
    just like $\QQ(\sqrt[3]2)$,
    although the two extensions are obviously not equal,
    since one of them only has real elements. In general,
    for an irreducible polynomial $p(x)\in\QQ[x]$,
    one can pick any of its roots to extend $\QQ$ by
    and still end up with an isomorphic extension,
    even if the various extensions are not equal.
    Thus, in some sense, the roots of a particular irreducible polynomial
    are algebraically indistinguishable, as they produce isomorphic
    field extensions.
\end{example}

We also have a converse to
Proposition \ref{prop:simple-extension-quotient-poly}:
if $[F(\alpha) : F] = n$ is finite,
then we can see that $\alpha$ must satisfy a polynomial equation because
the list $(1,\alpha,\dots,\alpha^n)$ has more than $n$ elements,
hence is $F$-linearly dependent.
By discarding factors of this polynomial we can force it to be irreducible,
and then the aforementioned proposition implies that its degree is $n$.
Thus we get the following statement:
\begin{prop}
    \label{prop:degree-of-primitive-element}
    If $F(\alpha)/F$ is a finite, simple extension,
    then there exists an irreducible polynomial with $\alpha$ as a root
    with degree equal to $[F(\alpha):F]$.
\end{prop}

\begin{exercise}
    If $F(\alpha)$ is a finite simple extension of $F$,
    show that there is a unique monic irreducible polynomial
    $m_{\alpha,F}(x)\in F[x]$ such that $m_{\alpha,F}(\alpha) = 0$.
    This polynomial is called the \vocab{minimal polynomial}
    of $\alpha$ over $F$.
\end{exercise}

\begin{prop}
    \label{prop:min-poly-divides-poly}
    Let $K/F$ be an extension.
    If $f(x)\in F[x]$ is a polynomial with root $\alpha\in K$, then
    $m_{\alpha,F}(x)$ divides $f(x)$.
\end{prop}

These propositions (and more to come) show that
$\alpha\in K$ being a root of a polynomial
with coefficients in $F\subseteq K$ 
gives us some nice information about the extension $F(\alpha)/F$.
In this vein, we define the following notion:
\begin{definition}[Algebraic extension, algebraic element]
    \label{def:algebraic}
    An element $\alpha$ of the field $K$ is
    \vocab{algebraic} over subfield $F$
    if $\alpha$ is the root of some polynomial in $F[x]$;
    equivalently (by Proposition \ref{prop:degree-of-primitive-element}),
    $F(\alpha)/F$ is a finite simple extension.
    The extension $K/F$ is an \vocab{algebraic extension}
    if each element of $K$ is algebraic over $F$.
\end{definition}

\begin{prop}
    \label{prop:finite-extensions-are-algebraic}
    If $K/F$ is a finite extension, it is algebraic.
\end{prop}

\begin{proof}
    If $K$ is finite-dimensional as an $F$-vector space,
    then, for any $\alpha\in K$,
    we see that $F(\alpha)\subseteq K$, being an $F$-linear subspace of $K$,
    is also finite-dimensional over $F$, so $\alpha$ is algebraic over $F$.
\end{proof}

This proposition essentially just uses the inequality
$[F(\alpha) : F] \le [K : F]$,
but can get a good amount more numerical information about degrees.

\begin{prop}[Tower law]
    \label{prop:tower-law}
    Let $K/F$ be a subextension of the finite extension $L/F$,
    i.e. $F\subseteq K\subseteq L$. Then
    \begin{equation*}
	[L : F] = [L : K][K : F]
    \end{equation*}
    where if one side is infinite then so is the other.
\end{prop}

\begin{proof}
    Let $\alpha_1$, \dots, $\alpha_n$ be an $F$-basis for $K$
    and $\beta_1$, \dots, $\beta_m$ be a $K$-basis for $L$.
    Then any $x\in L$ can be written as
    \begin{equation*}
        x = c_1\beta_1 + c_2\beta_2 + \cdots + c_m\beta_m
	= \sum_{i=1}^n\sum_{j=1}^m d_{i,j}\alpha_i\beta_j
    \end{equation*}
    for $c_i\in K$ and $d_{i,j}\in F$,
    so $\{\alpha_i\beta_j : 1\le i\le n,\; 1\le j\le m\}$ 
    is an $F$-spanning set for $L$. On the other hand if
    \begin{equation*}
	\sum_{i=1}^n\sum_{j=1}^m d_{i,j}\alpha_i\beta_j = 0
    \end{equation*}
    then the coefficients $\sum_{i=1}^n d_{i,j}\alpha_i$
    of the $\beta_j$ must each be zero
    by the linear independence of the $\beta_j$,
    and then the $d_{i,j}$ themselves are zero
    by the linear independence of the $\alpha_i$.
\end{proof}

An obvious corollary of this is that
for an extension $L/F$, any sub-extension $K/F$ must have
$[K : F]$ dividing $[L : F]$. Another example of
numerical information about degrees is the following:

\begin{prop}
    \label{prop:product-ineq}
    Let $K_1/F$ and $K_2/F$ be subextensions of the extension $L/F$
    and $K_1K_2$ be the smallest subfield of $L$ containing
    both $K_1$ and $K_2$. Then
    \begin{equation*}
	[K_1K_2 : F] \le [K_1 : F][K_2 : F].
    \end{equation*}
\end{prop}

\begin{proof}
    \todo{uhh just take product of bases?}
\end{proof}

\begin{prop}
    \label{prop:double-algebraic-extension}
    Let $L/K$ and $K/F$ be algebraic extensions.
    Then $L$ is algebraic over $F$.
\end{prop}

\begin{proof}
    Let $\alpha\in L$ and let its minimal polynomial over $K$ be
    $c_nx^n + \cdots + c_1x + c_0\in K[x]$.
    Then $\alpha$ is an algebraic element in
    $F(c_0,c_1,\dots,c_n)$, which is a finite extension of $F$
    because the $c_i$ are algebraic over $F$, so
    \begin{equation*}
	[F(\alpha) : F] \le [F(c_0,c_1,\dots,c_n,\alpha) : F] < \infty.
	\qedhere
    \end{equation*}
\end{proof}

\begin{remark}
    A non-algebraic extension $K/F$ is called \vocab{transcendental}.
    For any $t\in K$, the subring $F[t]$ is isomorphic to
    a quotient of the polynomial ring $F[x]$;
    if $t$ isn't algebraic over $F$ then we must have $F[t]\cong F[x]$,
    in which case $t$ is \vocab{transcendental} over $F$.
\end{remark}

We say that a field $\Omega$ is \vocab{algebraically closed}
if every polynomial in $\Omega$ has a root in $\Omega$,
or equivalently factors into a product of linear polynomials in $\Omega$,
or equivalently that the irreducible polynomials
in $\Omega$ all have degree $1$.
An \vocab{algebraic closure} of field $F$
is an algebraic extension $\Omega/F$ for which $\Omega$ is algebraically closed.

\begin{example}
    \label{ex:algebraic-closure-of-R}
    $\CC$ is algebraically closed (fundamental theorem of algebra),
    and it is an algebraic closure of $\RR$.
\end{example}

Important thing to know (not going to reproduce a proof here):
\begin{theorem}[Algebraic closures exist]
    \label{thm:algebraic-closures-exist}
    Every field $F$ has an algebraic closure,
    and any two algebraic closures of $F$ are isomorphic
    (though not at all uniquely).
\end{theorem}

\subsection{Splitting fields}
The key properties we will establish about field extensions
are all combinatorial in flavor, since finite-degree extensions
are fairly rigid objects.

\begin{prop}
    \label{prop:simple-extension-homomorphism-correspondence}
    Let $F(\alpha)/F$ be a simple extension,
    and $\varphi_0\colon F\to \Omega$ a homomorphism
    to some field $\Omega$.
    \begin{enumerate}[(a)]
        \ii If $\alpha$ is transcendental over $F$,
	there is a one-to-one correspondence
	between homomorphisms $F(\alpha)\to\Omega$ extending $\varphi_0$
	and transcendental elements of $\Omega$.
	\ii If $\alpha$ is algebraic over $F$
	with minimal polynomial $m(x)\in F[x]$,
	there is a one-to-one correspondence
	between homomorphisms $F(\alpha)\to\Omega$ extending $\varphi_0$
	and roots of $\varphi_0m$ in $\Omega$,
	where $\varphi_0m(x)\in F'[x]$ is the polynomial obtained
	by applying $\varphi_0$ to the coefficients of $m$.
    \end{enumerate}
\end{prop}

\begin{proof}
    \todo{pf}
\end{proof}

\begin{definition}[Splitting field]
    \label{def:splitting-field}
    Let $E/F$ be a field extension and $f(x)\in F[x]$.
    We say that \vocab{$f$ splits in $E$} or
    \vocab{$E$ splits $f$} if $f$ factors into
    a product of linear polynomials in $E[x]$, i.e.
    \begin{equation*}
        f(x) = c(x - \alpha_1)(x - \alpha_2)\cdots(x-\alpha_n),
	\quad c\in F,\; \alpha_i\in E.
    \end{equation*}
    A field $E$ in which $f$ splits is a \vocab{splitting field} for $f$
    if $E$ is also generated by the roots of $f$ in $E$.
\end{definition}

\begin{example}
    \label{ex:splitting-fields}
    (1) If $f(x) = ax^2 + bx + c\in\QQ[x]$ does not have rational roots,
    then $\QQ(\sqrt{b^2 - 4ac})$ is a splitting field for $f$.

    (2) Let $f(x)\in\QQ[x]$ be an irreducible cubic polynomial
    with roots $\alpha_1,\alpha_2,\alpha_3\in\CC$.
    Then $\QQ(\alpha_1,\alpha_2)\subseteq\CC$
    certainly is a splitting field for $f$
    (because we can recover $\alpha_3$ by Vieta's formulas).
    In $\QQ(\alpha_1)$ we can factor $f(x) = (x-\alpha_1)g(x)$
    for some degree $2$ polynomial $g(x)$, so,
    depending on the irreducibility of $g$ in $\QQ(\alpha_1)$,
    we get that the splitting field $\QQ(\alpha_1,\alpha_2)$
    either has degree $3$ or $6$ over $\QQ$ by the tower law.

    (3) Let $f(x) = 1 + x + \cdots + x^{p-1} = \frac{x^p - 1}{x-1}\in\QQ[x]$.
    If complex number $\zeta\ne 1$ satisfies $\zeta^p = 1$,
    then the roots of $f$ are $\zeta$, $\zeta^2$, \dots, $\zeta^{p-1}$,
    so $\QQ[\zeta]$ is a splitting field for $f$ with degree $p-1$.

    (4) Let $f(x) = x^p - x - a\in F[x]$,
    where $F$ has characteristic $p\ne 0$.
    Since $r^p = r$ for all $r$ in the prime subfield $\FF_p$,
    if $\alpha$ is a root of $f(x)$ in some extension $E$,
    then so are $\alpha+1$, \dots, $\alpha+p-1$,
    meaning $F(\alpha)$ is a splitting field for $f$.
\end{example}

In (2) of Example \ref{ex:splitting-fields},
we used the fact that $\CC$ is algebraically closed
to obtain a splitting field for $f$
(and we can do this in general for any polynomial in $\QQ[x]$).
In general, splitting fields always exist:
\begin{prop}
    \label{prop:splitting-fields-exist}
    Let $F$ be a field and $f(x)\in F[x]$ a polynomial.
    A splitting field $E/F$ of $f$ exists,
    and $[E : F]\le (\deg f)!$.
\end{prop}

\begin{proof}
    Let $f_0 = f$ and $F_0 = F$, and generate a sequence
    of polynomials and extensions of $F$ as follows:
    let $F_{n+1} = F_n(\alpha_n)$
    be a stem field for an irreducible factor of
    the polynomial $f_n$ (i.e. $F_{n+1}$ contains a root $\alpha_n$ of
    this factor and is generated by this root over $F_n$),
    and let $f_{n+1}(x) = f_n(x)/(x-\alpha_n)$.
    Clearly $E = F_{\deg f}$ is a splitting field for $f$.
    Since $\deg f_{n+1} = \deg f_n - 1$,
    $\deg f_n = \deg f - n$, so $[F_{n+1} : F_n]\le \deg f - n$,
    and the tower law says that
    \begin{equation*}
	[E : F] = \prod_{n=1}^{\deg f}[F_n : F_{n-1}]
	\le \prod_{n=1}^{\deg f} (\deg f - n) = (\deg f)!.\qedhere
    \end{equation*}
\end{proof}

The main question now about splitting fields
is how much leeway we have in their construction;
for example, to what extent do the particular sequences $F_n$ and $f_n$
used to generate a splitting field in
Proposition \ref{prop:splitting-fields-exist} change the field?

\begin{lemma}
    \label{lem:homomorphisms-into-splitting-field}
    Let $F$ be a field and $f(x)\in F[x]$.
    Let $E/F$ and $\Omega/F$ be extensions with
    $E$ generated by the roots of $f$ in $E$
    and $\Omega$ splitting $f$.
    There exists an $F$-homomorphism $\varphi\colon E\to\Omega$,
    and the number of such $F$-homomorphisms is bounded above by $[E : F]$,
    with equality if $f$ has distinct roots in $\Omega$.
\end{lemma}

\begin{proof}
    Let $\alpha_1$, \dots, $\alpha_m$ be the roots of $f$ in $E$
    and define $F_k = F(\alpha_1,\dots,\alpha_k)$
    for $1\le k\le m$ (so $F_m = E$), and $F_0 = F$.
    Furthermore, let $f_k$ be the minimal polynomial of $\alpha_{k+1}$
    in $F_k$ for $0\le k < m$, so $f_0 = f$.
    \todo{finish}
\end{proof}

\begin{corollary}
    \label{coro:splitting-fields-isomorphic}
    Let $F$ be a field and $f(x)\in F[x]$.
    Any two splitting fields $E/F$ and $E'/F$ are $F$-isomorphic.
\end{corollary}

\begin{proof}
    There exist $F$-homomorphisms $E\to E'$ and $E'\to E$,
    which respectively imply that $[E : F]\le[E' : F]$
    and $[E' : F]\le [E : F]$. Thus the two extensions
    have the same finite degree over $F$,
    so dimension counting implies that
    the $F$-homomorphisms are linear isomorphisms,
    which implies that they are bijective, hence also $F$-isomorphisms.
\end{proof}

Note that there is no canonical choice of isomorphism
between two splitting fields, so one still has to be careful
when talking about ``the splitting field'' of a polynomial.
One thing we can say about polynomials now is that
their factorizations in their splitting fields
are essentially unique. Suppose that $F$ is a field, $f(x)\in F[x]$,
and $E/F$ and $E'/F$ are splitting fields for $f$ in which
\begin{equation*}
    f(x) = c(x-\alpha_1)\cdots(x-\alpha_n)
    = c(x-\beta_1)\cdots(x-\beta_n),
    \quad c\in F,\; \alpha_i\in E,\;\beta_i\in E'
\end{equation*}
(note of course that the number of $\alpha_i$s and $\beta_i$s
is the same, namely equal to the degree of $f$).
Then, an $F$-isomorphism $\varphi\colon E\to E'$
induces an isomorphism of the respective polynomial rings
$\varphi\colon E[x]\to E'[x]$ by just sending $x$ to $x$, in which
\begin{equation*}
    \varphi f(x) = c\prod_{i=1}^n (x-\varphi(\alpha_i)),
\end{equation*}
so unique factorization implies that the $\varphi(\alpha_i)$
are a permutation of the $\beta_i$ (as multisets).
In particular, the multiplicities of the roots of a polynomial
in a splitting field are intrinsic to the polynomial
and independent of the choice of extension (as we should hope).

\begin{definition}[Separable polynomial]
    \label{def:separable-polynomial}
    Let $F$ be a field. A polynomial $f(x)\in F[x]$
    is \vocab{separable} if it has no multiple roots
    in any splitting field, and inseparable otherwise.
\end{definition}

There is an easy way to check whether a polynomial is separable.
Define the \vocab{formal derivative} of $f(x) = a_0 + a_1x + \cdots + a_nx^n$
to be the expected expression:
\begin{equation*}
    f'(x) = a_1 + 2a_2x + \cdots + na_nx^{n-1}.
\end{equation*}
This is a derivation on $F[x]$ (i.e. is an $F$-linear map
that satisfies the product rule).
The product rule implies that a root of $f(x)$ with multiplicity $m$
not divisible by the characteristic $p$ of $F$
(of course, this requires $p\ne 0$)
is also a root of $f'(x)$ with multiplicity $m-1$,
since if $f(x) = (x-\alpha)^m g(x)$ for some polynomial $g(x)$
with $g(\alpha)\ne 0$, then
\begin{equation*}
    f'(x) = (x-\alpha)^m g'(x) + m(x-\alpha)^{m-1}g(x)
    = (x-\alpha)^{m-1}((x-\alpha)g'(x) + mg(x))
\end{equation*}
and $(\alpha-\alpha)g'(\alpha) + mg(\alpha)$ is clearly nonzero at $0$
unless $p$ divides $m$.
If $m$ is divisible by $p$ then $f'(x) = (x-\alpha)^m g'(x)$,
so $\alpha$ is a root of $f'$ with multiplicity at least $m$.
Either way, if $f$ has a multiple root, then its derivative
also has that root (and you can keep track of its multiplicity,
which is especially nice in characteristic zero).
Conversely if $f$ does not have any multiple roots
then $m = 1$ above, which is not divisible
by any possible value of the characteristic.

\begin{prop}
    \label{prop:separability-criterion}
    A polynomial $f(x)\in F[x]$ is separable
    if and only if $f$ and $f'$ are relatively prime.
\end{prop}

This is fantastic, since derivatives are easy to compute
by just looking at the coefficients
and the greatest common divisor is easy to compute
via the Euclidean algorithm.

\begin{exercise}
    \label{exer:gcd-invariant-of-ambient-field}
    Let $f(x),g(x)\in F[x]$ and let $E/F$ be an extension.
    Show that the greatest common divisor of $f$ and $g$ in $F[x]$
    is the same as the greatest common divisor in $E[x]$.
    Thus, the greatest common divisor of $f$ and $f'$ is
    independent of the ambient field.
\end{exercise}

\begin{example}
    \label{ex:inseparable-polynomials}
    (1) Any irreducible polynomial over a characteristic zero field
    is separable, as are products of distinct irreducible polynomials.

    (2) Let $a$ be an element of the field $F$
    with characteristic $p$. The polynomial $x^p - x - a$ is separable,
    and it is irreducible if for example $a = 1$
    (see Example \ref{ex:splitting-fields}).

    (3) A field $F$ with characteristic $p\ne 0$ is \vocab{perfect}
    if $F^p \coloneq \{x^p : x\in F\}$ is equal to $F$.
    Then any irreducible polynomial in $F[x]$ is separable.
    If $f$ were inseparable, then $\gcd(f,f')\ne 1$
    and the irreducibility of $f$ implies that $f$ divides $f'$,
    which forces $f' = 0$ since $\deg f' < \deg f$.
    Thus $f(x) = g(x^p) = g(x)^p$ for some polynomial $g(x)\in F[x]$ because
    the coefficients of $g$ are $p$th powers, contradicting irreducibility.

    (4) If $F$ isn't a perfect field and $a\in F$ is not a $p$th power
    then $x^p - a$ is irreducible and inseparable,
    since if $\alpha^p = a$ then $x^p - a = (x-\alpha)^p$
    so $\alpha$ is the only root of the polynomial
    and has multiplicity $p$.
    In fact, this scenario, where the field is imperfect
    and the derivative is zero, is the only obstruction
    to an irreducible polynomial being inseparable.
\end{example}

\section{Automorphisms}
\begin{definition}[Automorphism]
    \label{def:automorphism}
    The group $\aut(E/F)$ is the group of
    $F$-automorphisms, which are $F$-isomorphisms $E\to E$.
\end{definition}

\begin{example}
    \label{ex:F-automorphisms}
    The conjugation map $\ol\bullet\colon\CC\to\CC$
    is an $\RR$-automorphism of $\CC$.
\end{example}

\begin{prop}
    \label{prop:degree-and-aut}
    
\end{prop}










\end{document}
